\documentclass{article}
\usepackage{amsmath,amssymb,amsthm}
\usepackage{algorithm,algorithmic}
\usepackage{bm}
\usepackage{enumitem}
\newtheorem{theorem}{Theorem}
\newtheorem{lemma}{Lemma}
\newtheorem{assumption}{Assumption}
\newcommand{\EE}{\mathbb{E}}
\newcommand{\RR}{\mathbb{R}}
\newcommand{\norm}[1]{\left\lVert #1\right\rVert}
\newcommand{\ip}[2]{\langle #1,#2\rangle}
\begin{document}

\section*{Randomized Coordinate Descent with Non‑Uniform Sampling (RCD‑NU)}

\section*{Mathematical Formulation}
Let $f:\RR^{d}\rightarrow\RR$ be a (possibly non‑convex) differentiable function.
For each coordinate $i\in\{1,\dots ,d\}$ we are given a sampling probability
\[
p_i>0,\qquad \sum_{i=1}^{d}p_i=1 .
\]
Denote by $e_i\in\RR^{d}$ the $i$‑th canonical basis vector and by
\[
g_i(w)\;:=\;\frac{\partial f}{\partial w_i}(w)
\]
the $i$‑th partial derivative.

\noindent\textbf{Update rule.}  
At iteration $t=0,1,\dots$ a coordinate $i_t$ is drawn independently
according to $\mathbb{P}(i_t=i)=p_i$. The iterate is updated only along that
coordinate:
\begin{equation}
\boxed{
w_{t+1}=w_t-\eta\, g_{i_t}(w_t)\,e_{i_t}},
\qquad\eta>0\ \text{(stepsize)} .
\label{eq:update}
\end{equation}
All other coordinates remain unchanged.

\section*{Pseudocode}
\begin{algorithm}[H]
\caption{Randomized Coordinate Descent with non‑uniform probabilities}
\label{alg:rcd}
\begin{algorithmic}[1]
\STATE \textbf{Input:} initial point $w_0\in\RR^{d}$, stepsize $\eta>0$, probabilities $\{p_i\}_{i=1}^{d}$
\FOR{$t=0,1,\dots,T-1$}
    \STATE Sample $i_t\sim\{1,\dots ,d\}$ with $\Pr(i_t=i)=p_i$
    \STATE Compute partial gradient $g_{i_t}(w_t)=\partial_{i_t}f(w_t)$
    \STATE $w_{t+1}=w_t-\eta\,g_{i_t}(w_t)\,e_{i_t}$ \hfill\COMMENT{only coordinate $i_t$ changes}
\ENDFOR
\STATE \textbf{Return} $w_T$
\end{algorithmic}
\end{algorithm}

\section*{Dry Run Example}
Consider the one‑dimensional quadratic
\[
f(w)=(w-3)^2,\qquad w_0=0,\ \eta=0.1 .
\]
Here $d=1$, $p_1=1$, $g_1(w)=2(w-3)$.

\begin{center}
\begin{tabular}{c|c|c|c}
$t$ & $w_t$ & $g_1(w_t)$ & $f(w_t)$\\\hline
0 & $0$ & $2(0-3)=-6$ & $9$\\
1 & $0-0.1(-6)=0.6$ & $2(0.6-3)=-4.8$ & $(0.6-3)^2=5.76$\\
2 & $0.6-0.1(-4.8)=1.08$ & $2(1.08-3)=-3.84$ & $(1.08-3)^2=3.6864$\\
3 & $1.08-0.1(-3.84)=1.464$ & $2(1.464-3)=-3.072$ & $(1.464-3)^2=2.3796$
\end{tabular}
\end{center}
The objective value decreases monotonically, illustrating the effect of the
coordinate update.

\newpage
\section*{Assumptions}
\begin{assumption}[Coordinate‑wise $L_i$‑smoothness]\label{as:smooth}
For each $i\in\{1,\dots ,d\}$ there exists $L_i>0$ such that for all $w\in\RR^{d}$ and
all $\alpha\in\RR$,
\[
\bigl|g_i(w+\alpha e_i)-g_i(w)\bigr|\le L_i|\alpha|.
\]
Equivalently,
\[
f(w+\alpha e_i)\le f(w)+\alpha\,g_i(w)+\frac{L_i}{2}\alpha^{2}.
\tag{A.1}
\]
\end{assumption}

\begin{assumption}[Bounded below]\label{as:bound}
There exists $f_{\inf}>-\infty$ such that $f(w)\ge f_{\inf}$ for all $w\in\RR^{d}$.
\end{assumption}

\begin{assumption}[Stepsize]\label{as:stepsize}
The stepsize satisfies
\[
0<\eta\le \frac{2}{L_{\max}},\qquad L_{\max}:=\max_{i}L_i .
\tag{A.2}
\]
\end{assumption}

No convexity is required; the results hold for general (possibly non‑convex)
functions satisfying the above smoothness and bounded‑below conditions.

\section*{Main Convergence Theorem}
\begin{theorem}\label{thm:main}
Let $\{w_t\}_{t\ge0}$ be generated by \eqref{eq:update} under
Assumptions \ref{as:smooth}–\ref{as:stepsize}.
Define the weighted gradient norm
\[
\| \nabla f(w)\|_{p^{-1}}^{2}
:=\sum_{i=1}^{d}\frac{g_i(w)^{2}}{p_i}.
\]
Then for any integer $T\ge1$,
\[
\frac{1}{T}\sum_{t=0}^{T-1}\EE\!\left[\| \nabla f(w_t)\|_{p^{-1}}^{2}\right]
\;\le\;
\frac{2\bigl(f(w_0)-f_{\inf}\bigr)}{\eta\,\bigl(2-L_{\max}\eta\bigr)}\;
\max_{i}\frac{1}{p_i}\;\frac{1}{T}.
\tag{T.1}
\]
Consequently,
\[
\min_{0\le t<T}\EE\!\left[\| \nabla f(w_t)\|_{p^{-1}}^{2}\right]
\;=\;O\!\left(\frac{1}{T}\right).
\]
\end{theorem}

\section*{Theoretical Properties}
\begin{itemize}[leftmargin=*]
\item \textbf{Stability.}  The condition $\eta\le 2/L_{\max}$ guarantees that the
coefficient $2-L_{\max}\eta$ in (T.1) is non‑negative, preventing divergence
even for non‑convex $f$.
\item \textbf{Hyperparameter sensitivity.}
A larger stepsize $\eta$ (up to the bound) improves the constant
$\eta(2-L_{\max}\eta)$, thus accelerating convergence; however, exceeding
$2/L_{\max}$ makes the right‑hand side of Lemma~\ref{lem:descent} negative,
invalidating the guarantee.
\item \textbf{Comparison to uniform sampling.}
If $p_i=1/d$ for all $i$, then $\max_i 1/p_i=d$ and (T.1) reduces to the classic
$O(d/T)$ bound for uniform randomized coordinate descent.
\item \textbf{Weighted norm interpretation.}
The factor $1/p_i$ up‑weights rarely sampled coordinates, showing that the
algorithm automatically compensates for the sampling bias in the convergence
rate.
\end{itemize}

\section*{Discussion}
The theorem shows that, despite the lack of convexity, the expected
weighted squared gradient norm decays at the optimal $1/T$ rate for
first‑order stochastic methods.  The appearance of $\max_i 1/p_i$ is inevitable
when sampling probabilities are non‑uniform: coordinates that are selected
infrequently dominate the convergence bound, which is precisely captured by
the weighted norm $\|\nabla f\|_{p^{-1}}$.

\newpage
\section*{Preliminaries (Definitions and Lemmas)}
\begin{lemma}[One‑step expected descent]\label{lem:descent}
Under Assumptions \ref{as:smooth}–\ref{as:stepsize},
\[
\EE\!\left[f(w_{t+1})\mid w_t\right]
\;\le\;
f(w_t)-\frac{\eta}{2}\bigl(2-L_{\max}\eta\bigr)
\sum_{i=1}^{d}p_i\,g_i(w_t)^{2}.
\tag{L.1}
\]
\end{lemma}
\begin{proof}
Fix $w_t$ and consider a particular coordinate $i$.
Using coordinate‑wise smoothness (A.1) with $\alpha=-\eta g_i(w_t)$ we obtain
\begin{align}
f\bigl(w_t-\eta g_i(w_t)e_i\bigr)
&\le f(w_t)-\eta g_i(w_t)^{2}
   +\frac{L_i}{2}\eta^{2}g_i(w_t)^{2} \nonumber\\
&= f(w_t)-\eta\Bigl(1-\frac{L_i\eta}{2}\Bigr)g_i(w_t)^{2}.
\tag{L.2}
\end{align}
Since $L_i\le L_{\max}$ and $\eta\le 2/L_{\max}$,
\[
1-\frac{L_i\eta}{2}\ge 1-\frac{L_{\max}\eta}{2}
=\frac{2-L_{\max}\eta}{2}\ge0 .
\tag{L.3}
\]
Taking expectation over the random choice of $i_t$,
\begin{align}
\EE\!\left[f(w_{t+1})\mid w_t\right]
&=\sum_{i=1}^{d}p_i\,f\bigl(w_t-\eta g_i(w_t)e_i\bigr) \nonumber\\
&\le\sum_{i=1}^{d}p_i\Bigl[f(w_t)-\eta\Bigl(1-\frac{L_i\eta}{2}\Bigr)g_i(w_t)^{2}\Bigr] \nonumber\\
&= f(w_t)-\eta\sum_{i=1}^{d}p_i\Bigl(1-\frac{L_i\eta}{2}\Bigr)g_i(w_t)^{2}.
\tag{L.4}
\end{align}
Using the uniform lower bound (L.3) for all $i$,
\[
\sum_{i=1}^{d}p_i\Bigl(1-\frac{L_i\eta}{2}\Bigr)g_i(w_t)^{2}
\ge \frac{2-L_{\max}\eta}{2}\sum_{i=1}^{d}p_i g_i(w_t)^{2}.
\tag{L.5}
\]
Substituting (L.5) into (L.4) yields exactly (L.1).  ∎
\end{proof}

\section*{Main Proof (Step‑by‑Step Derivation)}
We now prove Theorem~\ref{thm:main}.  All expectations are taken with respect to the
randomness in $\{i_0,\dots,i_{t-1}\}$ unless otherwise noted.

\begin{proof}
\noindent\textbf{[Step 1]}  Apply Lemma~\ref{lem:descent} and take total expectation:
\[
\EE\!\left[f(w_{t+1})\right]
\le \EE\!\left[f(w_t)\right]
-\frac{\eta}{2}\bigl(2-L_{\max}\eta\bigr)
\EE\!\left[\sum_{i=1}^{d}p_i\,g_i(w_t)^{2}\right].
\tag{P.1}
\]

\noindent\textbf{[Step 2]}  Rearrange (P.1) to isolate the weighted gradient term:
\[
\EE\!\left[\sum_{i=1}^{d}p_i\,g_i(w_t)^{2}\right]
\le \frac{2}{\eta\bigl(2-L_{\max}\eta\bigr)}
\Bigl(\EE\!\left[f(w_t)\right]-\EE\!\left[f(w_{t+1})\right]\Bigr).
\tag{P.2}
\]

\noindent\textbf{[Step 3]}  Sum (P.2) over $t=0,\dots,T-1$ and telescope the right‑hand side:
\begin{align}
\sum_{t=0}^{T-1}\EE\!\left[\sum_{i=1}^{d}p_i\,g_i(w_t)^{2}\right]
&\le \frac{2}{\eta\bigl(2-L_{\max}\eta\bigr)}
\Bigl(f(w_0)-\EE\!\left[f(w_T)\right]\Bigr) \nonumber\\
&\le \frac{2\bigl(f(w_0)-f_{\inf}\bigr)}{\eta\bigl(2-L_{\max}\eta\bigr)} .
\tag{P.3}
\end{align}
The last inequality uses Assumption~\ref{as:bound}.

\noindent\textbf{[Step 4]}  Relate the left‑hand side of (P.3) to the weighted norm
$\|\nabla f(w_t)\|_{p^{-1}}^{2}$.
For any non‑negative numbers $\{a_i\}$,
\[
\sum_{i=1}^{d}p_i a_i
= \sum_{i=1}^{d}\frac{a_i}{p_i}\,p_i^{2}
\ge \Bigl(\min_{i}p_i\Bigr)\sum_{i=1}^{d}a_i .
\tag{P.4}
\]
Applying (P.4) with $a_i=g_i(w_t)^{2}$ gives
\[
\sum_{i=1}^{d}p_i\,g_i(w_t)^{2}
\ge \bigl(\min_i p_i\bigr)\,\|\nabla f(w_t)\|^{2}.
\tag{P.5}
\]
Conversely,
\[
\|\nabla f(w_t)\|^{2}
= \sum_{i=1}^{d}g_i(w_t)^{2}
= \sum_{i=1}^{d}\frac{g_i(w_t)^{2}}{p_i}p_i
\le \Bigl(\max_i\frac{1}{p_i}\Bigr)
\sum_{i=1}^{d}p_i\,g_i(w_t)^{2}.
\tag{P.6}
\]

\noindent\textbf{[Step 5]}  Multiply both sides of (P.3) by $\max_i\frac{1}{p_i}$ and use (P.6):
\begin{align}
\sum_{t=0}^{T-1}\EE\!\left[\|\nabla f(w_t)\|_{p^{-1}}^{2}\right]
&= \sum_{t=0}^{T-1}\EE\!\left[\sum_{i=1}^{d}\frac{g_i(w_t)^{2}}{p_i}\right] \nonumber\\
&\le \Bigl(\max_{i}\frac{1}{p_i}\Bigr)
\sum_{t=0}^{T-1}\EE\!\left[\sum_{i=1}^{d}p_i\,g_i(w_t)^{2}\right] \nonumber\\
&\le \Bigl(\max_{i}\frac{1}{p_i}\Bigr)
\frac{2\bigl(f(w_0)-f_{\inf}\bigr)}{\eta\bigl(2-L_{\max}\eta\bigr)} .
\tag{P.7}
\end{align}

\noindent\textbf{[Step 6]}  Divide (P.7) by $T$ to obtain the average bound:
\[
\frac{1}{T}\sum_{t=0}^{T-1}\EE\!\left[\|\nabla f(w_t)\|_{p^{-1}}^{2}\right]
\le
\frac{2\bigl(f(w_0)-f_{\inf}\bigr)}{\eta\bigl(2-L_{\max}\eta\bigr)}
\max_{i}\frac{1}{p_i}\;\frac{1}{T},
\]
which is exactly the claim (T.1).  ∎
\end{proof}

\section*{Rate Derivation (With Constants)}
From Theorem~\ref{thm:main} we have the explicit constant
\[
C:=\frac{2\bigl(f(w_0)-f_{\inf}\bigr)}{\eta\bigl(2-L_{\max}\eta\bigr)}
\max_{i}\frac{1}{p_i},
\]
so that
\[
\frac{1}{T}\sum_{t=0}^{T-1}\EE\!\left[\|\nabla f(w_t)\|_{p^{-1}}^{2}\right]\le \frac{C}{T}.
\]
If one prefers a bound on the (unweighted) gradient norm, combine (P.5) and (P.3):
\[
\frac{1}{T}\sum_{t=0}^{T-1}\EE\!\left[\|\nabla f(w_t)\|^{2}\right]
\le
\frac{2\bigl(f(w_0)-f_{\inf}\bigr)}{\eta\bigl(2-L_{\max}\eta\bigr)}
\frac{1}{\min_i p_i}\;\frac{1}{T}.
\]

\section*{Special Cases}
\begin{enumerate}[leftmargin=*]
\item \textbf{Uniform sampling.}
If $p_i=1/d$ for all $i$, then $\max_i 1/p_i=d$ and the bound reduces to
\[
\frac{1}{T}\sum_{t=0}^{T-1}\EE\!\left[\|\nabla f(w_t)\|^{2}\right]
\le \frac{2d\bigl(f(w_0)-f_{\inf}\bigr)}{\eta\bigl(2-L_{\max}\eta\bigr)}\frac{1}{T},
\]
the classical $O(d/T)$ rate for uniform randomized coordinate descent.

\item \textbf{Deterministic cyclic coordinate descent.}
Setting $p_i=1/d$ and choosing the coordinates cyclically gives the same
expectation as uniform random sampling; the theorem therefore also applies as a
worst‑case bound for the cyclic scheme.

\item \textbf{Adaptive stepsize.}
If $\eta$ is chosen as $\eta_t=\frac{1}{L_{i_t}}$, then $2-L_{\max}\eta\ge1$ and the
constant improves to $\frac{2}{\eta_t}$, yielding a slightly better bound.
\end{enumerate}

\end{document}